\documentclass[../main.tex]{subfiles}

\begin{document}

\ifSubfilesClassLoaded{

    %\tableofcontents
    
    \setcounter{chapter}{1}
}{}

\chapter{ Ch2-4 , Ch2-5}
 代靖涵 25120222201319
\begin{exercise}{}{}
    设 $E \subset \mathbb{R}$, 且存在 $q: 0<q<1$, 使得对任一区间 $(a,b)$, 都有开区间列 $\{I_n\}$:
    \[
    E \cap (a,b) \subset \bigcup_{n=1}^{\infty} I_n, \quad \sum_{n=1}^{\infty} m(I_n) < (b-a)q,
    \]
    试证明 $m(E)=0$.
\end{exercise}
\begin{proof}
   存在 \(  q:0< q< 1  \), 使得 对于任意的 开区间\(  I  \), 都有 \[
   m^{*}\left( E\cap I \right)\le \sum _{n = 1}^{\infty}m\left( I_{n} \right)<   m\left( I \right)q  
   \] 若 \(  m^{*}\left( E \right)> 0   \),  则对于任意的 \(   \varepsilon > 0  \), 存在区间 \(  J  \), 使得 \[
   m\left( J \right)<  \left( 1+  \varepsilon  \right)m^{*}\left( J \cap E\right)   \implies m^{*}\left( E\cap J \right)>  \frac{1 }{1+  \varepsilon  }m\left( J \right)   
   \]  特别地, 取 \(   \varepsilon = \frac{1 }{q }-1> 0   \), 取 \(  I= J  \), 则导致一个矛盾.  因此 \(  m^{*}\left( E \right)= 0   \), 进而 \(  E  \)可测,  \(  m\left( E \right)= 0   \).   
\end{proof}
\begin{exercise}{}{}
    设 $\{B_\alpha\}_{\alpha \in I}$ 是 $\mathbf{R}^n$ 中的一族开球, 记 $G = \bigcup_{\alpha \in I} B_\alpha$. 若有 $0 < \lambda < m(G)$, 试证明存在有限个互不相交的开球 $B_{a_1}, B_{a_2}, \dots, B_{a_m}$, 使得

$$\sum_{k=1}^m m(B_{a_k}) > \frac{\lambda}{3^n}.$$
(作 $K \subset G$ 紧集: $m(K) > \lambda$, 再在 $\bigcup_{k=1}^m B_{a_k} \supset K$ 中取直径最大者, 后继者应与前者不相交. 最后直径放大三倍.)
\end{exercise}
\begin{proof}
    即我们需要证明对于任意的 \(  0< q< 1  \), 均有存在有限个互不相交的开球,使得 \[
    \sum _{k= 1}^{n}m\left( B_{\alpha _{k}} \right) > \frac{q }{3^{n} }m\left( \bigcup _{\alpha \in I}B_{\alpha } \right)  
    \] 由于 \(  G  \)可测, 则存在 闭集 \(  F  \)使得 \(  F\subseteq G  \), \(  m\left( F \right)> q m\left( G \right)    \),     \[
    \lim_{n\to \infty}m\left( F\cap \left( \bar{B_{n}} \cap \bar{G} \right) \right) = m\left( F \right)> qm\left( G \right)   
    \]其中 \(  B_{n}  \)以原点为中心, \(  n  \)为半径的开球.  故存在充分大的 \(  n_0  \), 使得 \[
    m\left( F\cap \bar{B}_{n_0}\cap  \bar{G} \right)> qm\left( G \right)  
    \]取 \(  K= F\cap  \bar{B}_{n_0} \cap  \bar{G} \), 则 \(  K  \)是一个有界闭集, 从而是紧的, 使得 \(  m\left( K \right)> q m\left( G \right)    \)  此外 \(  K\subseteq G  \). 因此    \(  \left\{ B_{\alpha } \right\}_{\alpha \in I}  \) 是 \(  K  \)的一个开覆盖, 存在有限的子覆盖 \(  B_{\beta  _1 },\cdots ,B_{\beta _{s}}  \). 不妨将 \(  B_{\beta _1 },\cdots ,B_{\beta _{s}}  \) 按直径从大到小排列, 将 \(  B_{\beta _1 }  \)和所有后续与 \(  B_{\beta _1 }  \)无交的成员取出, 构成 \[
    B_{\beta ^{2}_{1}},\cdots ,B_{\beta ^{2} _{s^{2}}}
    \]  将 \(  B_{\beta _1 ^{2}}  \)和 \(  B_{\beta _2 ^{2}}  \)取出, 并取出后续所有与 \(  B_{\beta _2 ^{2}}  \)无交的成员, 构成 \[
    B_{\beta _1 ^{3}}, B_{\beta _{2}^{3}},\cdots , B_{\beta ^{3}_{s^{3}}}
    \]   重复上述操作有限多次之后, 最后得到的成员记作 \[
    B_{\alpha _1 },\cdots ,B_{\alpha _{m}}
    \]  存在 某个 \( B_{\alpha _{i}} \), 使得   \(  B_{\beta _{j}}\cap B_{\alpha _{i}} \neq \varnothing  \),  \(  \mathrm{diam}\left( B_{\beta _{j}} \right)\le \mathrm{diam}\left(B_{\alpha _{i}} \right)    \), 则 \(  B_{\beta _{j}}\subseteq 3B_{\alpha _{i}}   \).   \(  3B_{\alpha _1 },\cdots ,3B_{\alpha _{m}}  \) 是 \(  K  \)的一个开覆盖, 从而 \[
    \sum _{k = 1}^{m}m\left( B_{\alpha _{k}} \right)= \frac{1 }{3^{n} }\sum _{k = 1}^{m}m\left( 3B_{\alpha _{k}} \right)\ge \frac{1 }{3^{n} }m\left( K \right)> \frac{q }{3^{n} }m\left( G \right)       
    \] 
\end{proof}
\begin{exercise}{}{}
设 $E \subset [0,1]$ 是可测集, 且有
\[m(E) \ge \varepsilon > 0, \quad x_i \in [0,1], \quad i=1,2,\cdots,n,\]
其中 $n > \frac{2}{\varepsilon}$, 试证明 $E$ 中存在两个点其距离等于 $\{x_1, x_2, \cdots, x_n\}$ 中某两个点之间的距离.
\end{exercise}
\begin{proof}
  定义 \(  E_{i}\coloneqq E-x_{i}  \), 则 \(  E_{i}  \)可测, 且 \(  m\left( E_{i} \right)= m\left( E \right)    \)   . 若 \(  E  \)中任意两点距离都不等于 \(  \left\{  x_1,\cdots,x_n  \right\}  \)中某两个点之间的距离, 则 断言\(  E_1,\cdots ,E_{n}  \)是两两无交的, 否则存在 \( k \in E_{i}\cap E_{j} \), 此时存在 \(y_{i},y_{j}\in E \), 使得  \(  y_{i}-x_{i}= y_{j}-x_{j}\implies y_{i}-y_{j}= x_{i}-x_{j}  \)  .   由于 \(  E\subseteq \left[ 0,1 \right], x_{i}\in \left[ 0,1 \right]    \), 我们有 \( E_{i}\subseteq \left[ -1,1 \right]   \), 但是 \[
  2= m\left( \left[ -1,1 \right]  \right)\ge m\left( \bigcup _{i= 1}^{n}E_{i} \right)  = \sum _{i = 1}^{n}m\left( E_{i} \right)\ge n \varepsilon > 2 
  \]矛盾.  
\end{proof}
\begin{exercise}{}{}
设 $W$ 是 $[0,1]$ 中的不可测集, 试证明存在 $\varepsilon: 0 < \varepsilon < 1$, 使得对于 $[0,1]$ 中任一满足 $m(E) \ge \varepsilon$ 的可测集 $E, W \cap E$ 是不可测集.
\end{exercise}

\begin{proof}
    若不然, 对于任意的 \(  n \in \mathbb{Z} _{> 0}  \), 存在可测集 \(  E_{n}\subseteq \left[ 0,1 \right]   \), 使得 \(  m\left( E_{n} \right)\ge 1-\frac{1 }{n }    \), 且 \(  W\cap E_{n}  \)可测.    令 \(  E\coloneqq \bigcup _{n = 1}^{\infty}E_{n}  \), 则 \(  E  \)  可测, 且 \[
    W\cap E= \bigcup _{n = 1}^{\infty}\left( W\cap E_{n} \right) 
    \]是可测的. 此时 \[
    \left( W\cap E \right)^{c}= W^{c}\cup E^{c} 
    \]可测. 其中 \[
    m\left( E^{c} \right)= m\left( \bigcap _{n = 1}^{\infty}E_{n}^{c} \right)\le m\left( E_{n}^{c} \right)= 1-m\left( E_{n} \right)= \frac{1 }{n }     , \forall n
    \]从而 \(  m\left( E^{c} \right)= 0   \). 进而 \(  m\left( E^{c}\setminus W^{c} \right)= 0   \), 我们有 \[
    W^{c}= \left( W^{c}\cup E^{c} \right)\setminus \left( E^{c}\setminus W^{c} \right)  
    \] 是可测的. 因而 \(  W  \)可测,   矛盾. 
\end{proof}
\end{document}